\hspace{3mm}

\centerline{\bf \huge 8 класс}

\begin{flushright}
        \scriptsize{ $-$ Если сдаваться, не попробовав, результат не изменится.}
\end{flushright}


\problem{1}
Расставьте по кругу четыре единицы, три двойки и три тройки так, чтобы сумма любых трёх подряд стоящих чисел не делилась на 3.

\problem{2}
На мехмате работают 88 преподавателей. Средняя зарплата преподавателей по мехмату равна 100 копеек. Средняя зарплата тех преподавателей, которые получают меньше 100 копеек, равна 85 копеек. Средняя зарплата тех преподавателей, которые получают больше 100 копеек, равна 135 копеек. Какое наименьшее число преподавателей могут иметь зарплату ровно 100 копеек? 

\problem{3}
Существуют ли натуральные числа $n$ и $k$, для которых $\dfrac{n}{3^k-n}$ --- квадрат натурального числа?

\problem{4}
В остроугольном треугольнике $A B C$ проведена высота $A H$ и отмечены середины $A_1, B_1$ и $C_1$ сторон $B C$, $CA$ и $A B$ соответственно. Точка $К$ симметрична точке $B_1$ относительно прямой $B C$. Докажите, что прямая $C_1 K$ делит отрезок $H A_1$ пополам.

\problem{5}
В 8 IT классе Техлицея некоторые школьники знакомы между собой, а некоторые нет. Оказалось, что если какой-то школьник знаком хотя бы с $20$ другими школьниками, то все его знакомые незнакомы между собой. Кроме того, если школьник незнаком хотя бы с 20 другими школьниками, то все его незнакомые знакомы между собой. Найдите максимальное возможное число школьников в таком классе.